\section{Conclusion}
\label{sec:conclusion}

We set out to predict the CCME Water Quality Index from physical and chemical
measurements and found that the conventional framing of that task does not survive
examination. The index is a deterministic function of the measurements that define it, so
a model given all eight recovers a formula rather than learning anything about water
quality. Measured directly, this framing yields $R^2 = 0.9997$ and a mean absolute error
of 0.100 index points -- figures that look like a result and are not one. We report them
only as an emulation benchmark and as a check that the pipeline is sound.

Reframing the problem as estimation under partial measurement availability restores its
content. Withholding orthophosphate alone moves XGBoost to $R^2 = 0.834$ and multiplies
error by a factor of 48; withholding half the measurements leaves $R^2 = 0.573$ and an
error of 8.596 index points, roughly the width of one CCME quality class. The step
between the full and partial regimes is the quantitative statement of how much apparent
accuracy was reconstruction rather than inference. Within this formulation XGBoost leads Random
Forest at every availability level, on identical test records and preprocessing, by a
margin that bootstrap intervals confirm is not a sampling artefact -- though the two were
trained on unequal budgets, so that margin bounds rather than measures the algorithmic
difference. Engineered
logarithmic and ratio features reduce error by 34\%, while flags marking upstream
imputation, being confounded with country, do not earn inclusion.

Two findings extend beyond this dataset. The first is that a published, reusable corpus
reporting no missing values had in fact been median-imputed before release, in a pattern
detectable only from the modal structure of its columns and concentrated almost entirely
in one contributing country. The second is that dissolved oxygen, which should be among
the most informative parameters in any water quality assessment, carries almost no
relationship to the index as distributed -- consistent with its guideline being applied in
the wrong direction upstream. Both were found by examining the data rather than the
documentation, and both would have been invisible to a study that reported only
predictive accuracy.

Phase 2 extends this work to a regularised multilayer perceptron trained on the same
partitions and preprocessing, which allows the tree ensembles and the neural model to be
compared without the confounds that make such comparisons unreliable in the literature.
The skewness documented in Section~\ref{sec:data} is expected to matter there in a way it
does not for trees, since neural models are not invariant to monotone transforms, and the
engineered features will be ablated accordingly. The point mass at the target's ceiling also
deserves direct treatment, through metrics separated by compliance status and potentially
through a model that represents the boundary explicitly rather than regressing through
it.
